\chapter{Généricité}
\sloppy

\section{Exercice 1 -- Classe Générique \texttt{Boite<T>}}

\subsection{Objectif}

Créer une classe générique \texttt{Boite<T>} avec les méthodes \texttt{setContenu}
et \texttt{getContenu}, testée avec \texttt{String} et \texttt{Integer}.

\subsection{Code source}

\begin{lstlisting}[language=Java, caption={exo1\_3.java -- Boite générique}]
public class exo1_3<T> {
    T contenu;

    void setContenu(exo1_3<T> e, T contenu) {
        e.contenu = contenu;
    }

    T getContenu(exo1_3<T> e) {
        return e.contenu;
    }

    public static void main(String[] args) {
        exo1_3<String> boiteS = new exo1_3<>();
        boiteS.setContenu(boiteS, "Java");
        System.out.println("Contenu String : " + boiteS.getContenu(boiteS));

        exo1_3<Integer> boiteI = new exo1_3<>();
        boiteI.setContenu(boiteI, 42);
        System.out.println("Contenu Integer : " + boiteI.getContenu(boiteI));
    }
}
\end{lstlisting}

\subsection{Explication}

Le paramètre de type \texttt{T} est résolu à la compilation selon le type
utilisé lors de l'instanciation. Cela garantit la \textbf{sûreté de type}
(\textit{type-safety}) sans duplication de code.

\section{Exercice 2 -- Classe \texttt{Paire<T, U>}}

\subsection{Objectif}

Créer une classe générique à deux paramètres de type avec une méthode
\texttt{afficherPaire()}.

\subsection{Code source}

\begin{lstlisting}[language=Java, caption={exo2\_3.java -- Paire générique}]
public class exo2_3<T, U> {
    T premier;
    U second;

    public exo2_3(T premier, U second) {
        this.premier = premier;
        this.second = second;
    }

    void afficherPaire(exo2_3<T, U> p) {
        System.out.println("Paire -> Premier: " + p.premier
                + ", Second: " + p.second);
    }

    public static void main(String[] args) {
        exo2_3<String, Integer> paire = new exo2_3<>("Java", 2026);
        paire.afficherPaire(paire);
    }
}
\end{lstlisting}

\section{Exercice 3 -- Interface Générique \texttt{Calcul<T>}}

\subsection{Objectif}

Définir une interface générique \texttt{Calcul<T>} et l'implémenter pour
\texttt{Integer} et \texttt{Double}.

\subsection{Code source}

\begin{lstlisting}[language=Java, caption={exo3\_3.java -- Interface Calcul générique}]
interface Calcul<T> {
    T addition(T a, T b);
}

class CalculInt implements Calcul<Integer> {
    public Integer addition(Integer a, Integer b) {
        return a + b;
    }
}

class CalculDouble implements Calcul<Double> {
    public Double addition(Double a, Double b) {
        return a + b;
    }
}
\end{lstlisting}

\section{Exercice 4 -- Interface \texttt{Comparateur<T>}}

\subsection{Objectif}

Définir une interface \texttt{Comparateur<T>} et l'implémenter pour
\texttt{Integer} (par valeur) et \texttt{String} (par longueur).

\subsection{Code source}

\begin{lstlisting}[language=Java, caption={exo4\_3.java -- Interface Comparateur}]
interface Comparateur4<T> {
    int comparer(T a, T b);
}

class ComparateurInt4 implements Comparateur4<Integer> {
    public int comparer(Integer a, Integer b) {
        return a.compareTo(b);
    }
}

class ComparateurString4 implements Comparateur4<String> {
    public int comparer(String a, String b) {
        return Integer.compare(a.length(), b.length());
    }
}
\end{lstlisting}

\subsection{Explication}

La valeur retournée par \texttt{comparer()} suit la convention Java~:
négatif si $a < b$, zéro si $a = b$, positif si $a > b$.

\section{Exercice 5 -- Deuxième Implémentation de \texttt{Comparateur<T>}}

\subsection{Objectif}

Réimplémenter l'interface \texttt{Comparateur<T>} dans un fichier indépendant
pour consolider la compréhension.

\subsection{Code source}

\begin{lstlisting}[language=Java, caption={exo5\_3.java -- Comparateur5}]
interface Comparateur5<T> {
    int comparer(T a, T b);
}

class ComparateurInt5 implements Comparateur5<Integer> {
    public int comparer(Integer a, Integer b) {
        return a.compareTo(b);
    }
}

class ComparateurString5 implements Comparateur5<String> {
    public int comparer(String a, String b) {
        return Integer.compare(a.length(), b.length());
    }
}
\end{lstlisting}

\section{Exercice 6 -- Méthodes Génériques Statiques}

\subsection{Objectif}

Créer deux méthodes statiques génériques~: \texttt{afficherTableau(T[])}
et \texttt{getPremier(T[])}.

\subsection{Code source}

\begin{lstlisting}[language=Java, caption={exo6\_3.java -- Méthodes génériques statiques}]
public class exo6_3 {

    public static <T> void afficherTableau(T[] tableau) {
        for (T element : tableau) {
            System.out.print(element + " ");
        }
        System.out.println();
    }

    public static <T> T getPremier(T[] tableau) {
        if (tableau.length > 0) {
            return tableau[0];
        }
        return null;
    }

    public static void main(String[] args) {
        Integer[] nombres = {10, 20, 30, 40, 50};
        String[] mots = {"Java", "Data", "ENSAH"};

        afficherTableau(nombres);
        System.out.println("Premier : " + getPremier(nombres));
        afficherTableau(mots);
        System.out.println("Premier : " + getPremier(mots));
    }
}
\end{lstlisting}

\section{Exercice 7 -- Borne Supérieure \texttt{<T extends Number>}}

\subsection{Objectif}

Créer une méthode générique bornée qui calcule la somme de deux nombres
de n'importe quel type numérique.

\subsection{Code source}

\begin{lstlisting}[language=Java, caption={exo7\_3.java -- Borne supérieure}]
public class exo7_3 {

    public static <T extends Number> double somme(T a, T b) {
        return a.doubleValue() + b.doubleValue();
    }

    public static void main(String[] args) {
        System.out.println("Somme : " + somme(5, 3.14));
    }
}
\end{lstlisting}

\subsection{Explication}

\begin{accentbox}
    \centering\vspace{0.3cm}
    {\large La borne \texttt{<T extends Number>} restreint le paramètre
    de type aux sous-classes de \texttt{Number} (\texttt{Integer},
    \texttt{Double}, \texttt{Float}\ldots), permettant d'appeler
    \texttt{doubleValue()} en toute sécurité.}
    \vspace{0.3cm}
\end{accentbox}

\section{Exercice 8 -- Héritage Générique (Animal / Chien)}

\subsection{Objectif}

Créer une classe générique \texttt{Animal<T>} et une sous-classe
\texttt{Chien} qui fixe le type à \texttt{String}.

\subsection{Code source}

\begin{lstlisting}[language=Java, caption={exo8\_3.java -- Héritage générique}]
class Animal<T> {
    T nom;
}

class Chien extends Animal<String> {
    public Chien(String nom) {
        this.nom = nom;
    }
}

public class exo8_3 {
    public static void main(String[] args) {
        Chien c = new Chien("Rex");
        System.out.println("Nom du chien : " + c.nom);
    }
}
\end{lstlisting}

\section{Exercice 9 -- Héritage Paramétré (Vehicule / Voiture)}

\subsection{Objectif}

Créer une classe générique \texttt{Vehicule<T>} et une sous-classe
\texttt{Voiture<T>} qui hérite du paramètre de type.

\subsection{Code source}

\begin{lstlisting}[language=Java, caption={exo9\_3.java -- Vehicule et Voiture génériques}]
class Vehicule<T> {
    T vitesse;
}

class Voiture<T> extends Vehicule<T> {
    public Voiture(T vitesse) {
        this.vitesse = vitesse;
    }
}

public class exo9_3 {
    public static void main(String[] args) {
        Voiture<Integer> v = new Voiture<>(120);
        System.out.println("Vitesse : " + v.vitesse);
    }
}
\end{lstlisting}

\section{Exercice 10 -- Repository Générique}

\subsection{Objectif}

Créer un \texttt{Repository<T>} générique et une classe
\texttt{UserRepository} qui le spécialise pour les objets \texttt{User}.

\subsection{Code source}

\begin{lstlisting}[language=Java, caption={exo10\_3.java -- Repository générique}]
class User {
    String nom;
    public User(String nom) { this.nom = nom; }
    public String toString() { return "User: " + nom; }
}

class Repository<T> {
    void save(T obj) {
        System.out.println("Objet sauvegarde : " + obj.toString());
    }
}

class UserRepository extends Repository<User> {}

public class exo10_3 {
    public static void main(String[] args) {
        User u = new User("Yazid");
        UserRepository repo = new UserRepository();
        repo.save(u);
    }
}
\end{lstlisting}

\section{Exercice 11 -- Wildcard \texttt{<?>}}

\subsection{Objectif}

Créer une méthode \texttt{afficherListe(List<?> liste)} acceptant
n'importe quel type de liste.

\subsection{Code source}

\begin{lstlisting}[language=Java, caption={exo11\_3.java -- Wildcard ?}]
import java.util.List;
import java.util.ArrayList;

public class exo11_3 {

    void afficherListe(List<?> liste) {
        for (Object obj : liste) {
            System.out.print(obj + " ");
        }
        System.out.println();
    }

    public static void main(String[] args) {
        List<String> noms = new ArrayList<>();
        noms.add("Yazid");
        noms.add("ENSAH");

        exo11_3 e = new exo11_3();
        e.afficherListe(noms);
    }
}
\end{lstlisting}

\subsection{Explication}

Le wildcard \texttt{<?>} signifie \og n'importe quel type \fg{}. Il permet
d'écrire une méthode générique sans spécifier de paramètre de type, mais
en contrepartie on ne peut \textbf{pas} ajouter d'éléments à la liste
(sauf \texttt{null}).

\section{Exercice 12 -- Wildcard Borné \texttt{<? extends Number>}}

\subsection{Objectif}

Créer une méthode \texttt{afficherNombres(List<? extends Number>)} testée
avec \texttt{List<Integer>} et \texttt{List<Double>}.

\subsection{Code source}

\begin{lstlisting}[language=Java, caption={exo12\_3.java -- Wildcard borné}]
import java.util.List;
import java.util.ArrayList;

public class exo12_3 {

    void afficherNombres(List<? extends Number> liste) {
        for (Number n : liste) {
            System.out.print(n + " ");
        }
        System.out.println();
    }

    public static void main(String[] args) {
        List<Integer> entiers = new ArrayList<>();
        entiers.add(10);
        entiers.add(20);

        List<Double> reels = new ArrayList<>();
        reels.add(15.5);
        reels.add(25.5);

        exo12_3 e = new exo12_3();
        e.afficherNombres(entiers);
        e.afficherNombres(reels);
    }
}
\end{lstlisting}

\subsection{Comparaison des Wildcards}

\begin{table}[H]
\centering
\begin{tabularx}{0.95\textwidth}{C{3.5cm} X X}
\toprule
\textbf{\color{primarycolor}Wildcard} &
\textbf{\color{primarycolor}Signification} &
\textbf{\color{primarycolor}Usage typique} \\
\midrule
\texttt{<?>} & N'importe quel type & Lecture seule, type inconnu \\
\addlinespace[0.15cm]
\texttt{<? extends T>} & T ou ses sous-types & Lecture d'une collection de T \\
\addlinespace[0.15cm]
\texttt{<? super T>} & T ou ses super-types & Écriture dans une collection \\
\bottomrule
\end{tabularx}
\caption{Comparaison des wildcards en Java}
\end{table}
